\documentclass[11pt]{letter}
\usepackage[margin=1in]{geometry}
\signature{Arjun Pemmasani\ on behalf of both authors}
\address{Harvey Mudd College\ Claremont, CA 91711\ apemmasani@hmc.edu}
\begin{document}
\begin{letter}{The Editor\ American Journal of Physics}
\opening{Dear Editor,}

We submit ``How long is a simulated double pendulum right? Step size,
floating-point precision, and the Lyapunov exponent'' for consideration in the
Computational Physics section.

The manuscript is a teaching contribution and presents no new physics. It
describes a computational exercise for upper-division courses in which students
integrate the double pendulum in arbitrary-precision arithmetic, vary the step
and the precision independently, and find that the time for which the simulation
is accurate obeys one logarithmic law governed by the largest Lyapunov exponent.
Students measure that exponent from their own numerical error and compare it with
the standard two-trajectory estimate. Along the way they meet a best step for each
floating-point format, a ceiling on what double precision can achieve, and a
trajectory that is entirely wrong while conserving energy to many digits. The
manuscript ends with six suggested problems.

The Journal has treated the double pendulum as a vehicle for teaching chaos since
Shinbrot \emph{et al.}\ [\textbf{60}, 491 (1992)] and Levien and Tan
[\textbf{61}, 1038 (1993)]. Those papers concern the physical system and how it
amplifies an uncertainty in its initial state. Ours concerns the simulation that
instructors and students now routinely substitute for the apparatus, and asks how
long it can be believed. We are not aware of a treatment of this question at the
undergraduate level.

In accordance with the Journal's procedure the manuscript file contains no author
names, affiliations or acknowledgments. The supplementary material comprises the
source code, the plotted data and two animations. The work is original and is not
under consideration elsewhere, and both authors have approved its submission.

\closing{Sincerely,}
\end{letter}
\end{document}
